%----------
%	CONTENIDO NEGATIVO - MACHINE LEARNING
%----------	


We have followed a similar approach as in Section~\ref{sec:ageneral_ml}, with a total of 158 experiments in the span of five months (24/02/24 - 09/07/24), and defining five different windows illustrated in Figures ~\ref{fig:exp_cnegativo} and ~\ref{fig:exp_cnegativo_classes}. Figure \ref{fig:exp_cnegativo} shows the overall trend in weighted average F1 scores, and we can see that the global trend line, depicted in red, with an $R^2 = 0.13$, highlights a moderate yet positive impact of the changes made in each window. Figure \ref{fig:exp_cnegativo_classes}, on the other hand, breaks down the performance for specific classes: "Insultos", "Desprestigiar Acto", "Desprestigiar Víctima", and "Desprestigiar Deportista Autora". Each window includes its own trend line, offering a more granular view of how F1-Scores evolve and respond to various modifications over time. Notably, "Desprestigiar Deportista Autora" consistently demonstrated high performance, with F1 scores that remained above the overall trend line. In contrast, the majority of experiments for the other classes tend to fall below this trend.


\begin{figure}[H]
	\ffigbox[\FBwidth]
	{\caption{Distribution of the ML Experiments over Time for "Contenido Negativo" with the Different Windows and Weighted Avg. F1-Score Trend Line}\label{fig:exp_cnegativo}}
 % [Aquí ponemos como queremos que aparezca en el índice, sin la referencia]{Aquí como queremos que aparezca en el documento, con la referencia}
	{\includegraphics[scale=0.25]{Plantilla_TFG_ingles_2019/imagenes/Experiments in _ContenidoNegativo_ over time.pdf}}
\end{figure}

\begin{figure}[H]
	\ffigbox[\FBwidth]
	{\caption{Distribution of the ML Experiments over Time for "Contenido Negativo", per Class and with the Different Windows (Trend Line of Weighted Avg. F1-Score in Red)}\label{fig:exp_cnegativo_classes}}
 % [Aquí ponemos como queremos que aparezca en el índice, sin la referencia]{Aquí como queremos que aparezca en el documento, con la referencia}
	{\includegraphics[scale=0.25]{Plantilla_TFG_ingles_2019/imagenes/Experiments in _ContenidoNegativo_ over time_ perclass_ trendline.pdf}}
\end{figure}


\begin{table}[H]
    \centering
    \begingroup
    \small % Change to desired font size
    \ttabbox[\FBwidth]{
        \caption{CNEG Average Evaluation Metrics per Window in ML Experiments}
        \label{tab:window-metrics-cnegativo} % Unique label for the table
    }{
        \begin{tabularx}{\textwidth}{
            X
            >{\centering\arraybackslash}X
            >{\centering\arraybackslash}X
            >{\centering\arraybackslash}X
            >{\centering\arraybackslash}X
            >{\columncolor{gray!15}\centering\arraybackslash}X % Highlight this column
            >{\centering\arraybackslash}X
            >{\centering\arraybackslash}X
            >{\centering\arraybackslash}X
            >{\centering\arraybackslash}X
        }
            \hline
            \rowcolor{gray!30} % Adding a gray color to the header row
            \textbf{} & \textbf{Acc. (Std.)} & \textbf{Time (s)} & \textbf{W. Avg. Prec.} & \textbf{W. Avg. Rec.} & \textbf{W. Avg. F1-Score} & \textbf{D. Víc. F1-Score} & \textbf{D. Acto F1-Score} & \textbf{Ins. F1-Score} & \textbf{D. D. Aut. F1-Score} \\
            \hline
            W1 & 0.486 (0.041) & 126.235 & 0.492 & 0.486 & 0.482 & 0.438 & 0.292 & 0.325 & 0.648 \\
            \hline
            W2 & 0.476 (0.083) & 1151.250 & 0.484 & 0.476 & 0.475 & 0.389 & 0.349 & 0.425 & 0.595 \\
            \hline
            W3 & 0.468 (0.054) & 1088.658 & 0.480 & 0.468 & 0.466 & 0.365 & 0.362 & 0.421 & 0.584 \\
            \hline
            W4 & 0.484 (0.051) & 757.250 & 0.493 & 0.484 & 0.465 & 0.337 & 0.347 & 0.381 & 0.618 \\
            \hline
            W5 & 0.570 (0.058) & 3023.991 & 0.588 & 0.570 & 0.574 & 0.640 & 0.381 & 0.429 & 0.697 \\
            \hline
            \multicolumn{10}{l}{Source: 'Training.xlsx'} \\
            \hline
        \end{tabularx}
    }
    \endgroup
\end{table}




\begin{itemize}
    \item \textbf{W1: Baseline Establishment:} In this initial phase, the dataset remained unmodified, employing FastText embedding (as detailed in Section~\ref{text_encoding_ml}) with Random Forest and Logistic Regression models. While the mean performance in terms of weighted F1-Score was the second highest with 0.482, "Desprestigiar Acto" achieved the lowest performance of the experiments with 0.292, as seen in Table \ref{tab:window-metrics-cnegativo}.


    \item \textbf{W2: Advanced Features, Embeddings and Models Integration:} During this period, several significant changes were introduced:
        \begin{itemize}
            \item \textbf{Data:} Stylistic features were added as described in Section~\ref{stylistic_features}, the 'Comentario Neutro' category was removed, and comprehensive text preprocessing was applied as detailed in Section~\ref{sec:data_preprocessing}.
            \item \textbf{Balancing:} SMOTE was employed to handle data imbalance.
            \item \textbf{Embedding:} Advanced embeddings detailed in Section~\ref{text_encoding_dl} were incorporated into the analysis.
            \item \textbf{Models:} The rest of the traditional ML models detailed in Section~\ref{sec:ml_models} were introduced (MLP, XGBoosting, naive Bayes)
        \end{itemize}
    From Table \ref{tab:window-metrics-cnegativo} we can see that the F1-scores during this window showed a varied response to these modifications, with improvements in F1-Score for "Desprestigiar Acto" and "Insultos". Additionally, although the computation time increased by 9.119, it did not imply an increase in the general performance.


    \item \textbf{W3: ADASYN Implementation:} The settings for this phase were equal to those of W2 but testing such modifications with ADASYN. In this case, in Table \ref{tab:window-metrics-cnegativo} the performance metrics during this window revealed that F1-Score decreased for all classes, unless for "Desprestigiar Acto" (0.349 vs 0.362). This suggested that ADASYN may not be the most effective data-balancing method for this specific task.


    \item \textbf{W4: Control Phase Without Balancing Techniques:} This window served as a control phase where all modifications were tested without applying any balancing techniques. From Figure \ref{fig:exp_cnegativo_classes} we can see that trend lines have a negative slope in this window, hence, resulting in a decline in model performance, as the only class that improved was "Desprestigiar Deportista Autora" with F1-Score of 0.618.


    \item \textbf{W5: Introduction of OpenAI Embedding:} In the final phase, OpenAI "text-embedding-3-large" embedding was introduced, leading to the highest F1-scores observed throughout the experimentation in all classes. The overall weighted F1-Score reached 0.574 and 0.697 for "Desprestigiar Deportista Autora". Additionally, in Figure~\ref{fig:exp_cnegativo_classes} we can see that only the experiments for "Desprestigiar Deportista Autora" and "Desprestigiar Víctima" were above the trend line, suggesting some room for improvements for "Desprestigiar Acto" and "Insultos".


\end{itemize}


From this analysis, we provide the following tables with the top 5 best overall models for \textit{Contenido Negativo}. Table \ref{tab:cnegativo-bestmodels} lists the models and their corresponding parameters. Meanwhile, Table \ref{tab:cnegativo-bestmodels_metrics} shows the performance metrics for these models.

\begin{table}[H]
    \centering
    \begingroup
    \small % Change to desired font size
    \ttabbox[\FBwidth]{
        \caption{CNEG Best ML Models Nomenclature}
        \label{tab:cnegativo-bestmodels} % Unique label for the table
    }{
        \begin{tabularx}{\textwidth}{
            X  % Name
            X  % Model
            X  % Embedding
            X  % CV
            X  % Balance
            }
            \hline
            \rowcolor{gray!30} % Adding a gray color to the header row
            \textbf{} & \textbf{Model} & \textbf{Embedding} & \textbf{CV} & \textbf{Balance} \\
            \hline
            \multicolumn{5}{l}{\textbf{Naive Models}} \\
            \hline
            \cnegmajor & - & - & - & - & - \\
            \hline
            \cnegrandom & - & - & - & - & - \\
            \hline
            \multicolumn{5}{l}{\textbf{Finetuned Models}} \\
            \hline
            \cnegmlptead & MLP & text-embedding-3-large & 5 & ADASYN \\
            \hline
            \cneglrteno & Logistic Regression & text-embedding-3-large & 5 & None \\
            \hline
            \cnegxgbtead & XGBoosting & text-embedding-3-large & 5 & ADASYN \\
            \hline
            \cnegrftead & Random Forest & text-embedding-3-large & 5 & ADASYN \\
            \hline
            \cnegxgbteno & XGBoosting & text-embedding-3-large & 5 & None \\
            \hline
            \multicolumn{5}{l}{Source: 'Training.xlsx'} \\
            \hline
        \end{tabularx}
    }
    \endgroup
\end{table}


\begin{table}[H]
    \centering
    \begingroup
    \small % Change to desired font size
    \ttabbox[\FBwidth]{
        \caption{CNEG Evaluation Metrics for the Best ML Models}
        \label{tab:cnegativo-bestmodels_metrics} % Unique label for the table
    }{
        \begin{tabularx}{\textwidth}{
            X
            >{\centering\arraybackslash}X
            >{\columncolor{gray!15}\centering\arraybackslash}X % Highlight this column
            >{\centering\arraybackslash}X
            >{\centering\arraybackslash}X
            >{\centering\arraybackslash}X
            >{\centering\arraybackslash}X
        }
            \hline
            \rowcolor{gray!30} % Adding a gray color to the header row
            \textbf{} & \textbf{Acc. (Std.)} & \textbf{W. Avg. F1-Score} & \textbf{D. Víc. F1-Score} & \textbf{D. Acto F1-Score} & \textbf{Ins. F1-Score} & \textbf{D. D. Aut. F1-Score} \\
            \hline
            \multicolumn{7}{l}{\textbf{Naive Models}} \\
            \hline
            \cnegmajor & 0.421 (0.027) & 0.250 & 0.000 & 0.000 & 0.000 & 0.592 \\
            \hline
            \cnegrandom & 0.289 (0.024) & 0.289 & 0.201 & 0.205 & 0.169 & 0.420 \\
            \hline
            \multicolumn{7}{l}{\textbf{Finetuned Models}} \\
            \hline
            \cnegmlptead & 0.634 (0.064) & 0.630 & 0.645 & 0.483 & 0.455 & 0.767 \\
            \hline
            \cneglrteno & 0.620 (0.071) & 0.620 & 0.690 & 0.467 & 0.417 & 0.746 \\
            \hline
            \cnegxgbtead & 0.592 (0.067) & 0.602 & 0.643 & 0.457 & 0.480 & 0.704 \\
            \hline
            \cnegrftead & 0.592 (0.065) & 0.601 & 0.741 & 0.400 & 0.435 & 0.702 \\
            \hline
            \cnegxgbteno & 0.592 (0.059) & 0.594 & 0.643 & 0.424 & 0.455 & 0.712 \\
            \hline
            \multicolumn{7}{l}{Source: 'Training.xlsx'} \\
            \hline
        \end{tabularx}
    }
    \endgroup
\end{table}


From this analysis, we have decided to choose model \textbf{\cnegmlptead} as the best-performing model for this task, as it achieved the highest weighted average F1-Score of 0.630 while maintaining high performance across other classes with F1-Scores of 0.645, 0.483, 0.455 and 0.767 for "Desprestigiar Víctima", "Desprestigiar Acto", "Insultos" and "Desprestigiar Deportista Autora", respectively. Moreover, it achieved the second lowest standard deviation of the models detailed in Table~\ref{tab:cnegativo-bestmodels_metrics}, and significant improvements compared to naive models {\cnegmajor} and {\cnegrandom}.

We provide the confusion matrix for this model in Figure \ref{fig:exp_cnegativo_metrics}. In particular, we can see that overall the model performed consistently in all classes, particularly at correctly classifying "Desprestigiar Deportista Autora", with 23 correct predictions, but struggled significantly with other categories, particularly "Desprestigiar Víctima", "Desprestigiar Acto" and "Insultos". One possible reason behind this is that the context and meaning of such classes do not differ much between them, hence the misclassification.


\begin{figure}[H]
	\ffigbox[\FBwidth]
	{\caption{Confusion Matrix for the Best ML Model (\textbf{\cnegmlptead}) for "Contenido Negativo"}\label{fig:exp_cnegativo_metrics}}
 % [Aquí ponemos como queremos que aparezca en el índice, sin la referencia]{Aquí como queremos que aparezca en el documento, con la referencia}
	{\includegraphics[scale=0.11]{Plantilla_TFG_ingles_2019/imagenes/Experiments in _ContenidoNegativo_ best_model.pdf}}
\end{figure}
